\section{1.1 Background and Context}
Intelligent sports training systems are moving from fixed mechanical routines toward adaptive decision-making. In football-specific training, launcher machines can generate repeatable balls but commonly operate as programmable throwers, not as situationally aware systems. In most real deployments, operator experience and manual setup still dominate targeting quality. This practical gap becomes critical when training objectives require consistent stimulation of specific body parts, directional responses, or tactical patterns.

The target use case of this thesis is a body-part-targeted football training scenario where the athlete occupies a constrained indoor arena and the machine should eventually interpret commands such as \texttt{left shoulder}, \texttt{right knee}, \texttt{body center}, or \texttt{goal zone A}. The core requirement for this vision is a unified metric 3D understanding of the scene: camera poses, player keypoints, ball trajectory, and virtual/physical targets must coexist in one world frame.

This requirement reframes the problem from "detection" to "geometric consistency and decision readiness". A visually appealing 2D overlay is insufficient if triangulated 3D points drift by 20-40 cm. For launcher control, these errors directly map to dangerous or ineffective shots. Therefore, this thesis focuses on building a measurable perception backbone with explicit error budgets and operational constraints.

\section{1.2 Problem Statement}
The research problem addressed is:

How can a practical four-camera vision system be calibrated, synchronized, validated, and integrated so that it provides reliable 3D positions of ball and body keypoints in a fixed indoor arena with accuracy sufficient for staged closed-loop launcher control?

This problem includes five tightly coupled sub-problems:
\begin{enumerate}
  \item Camera intrinsics must be stable and geometry-faithful at target runtime resolution.
  \item Extrinsics must remain valid despite physical adjustments and partial marker visibility.
  \item Multi-camera synchronization and triangulation must tolerate real-world frame inconsistencies.
  \item Ball and human 2D detections must be converted to robust 3D with controllable outliers.
  \item Error must be quantified against protocolized GT before any control-loop claim is made.
\end{enumerate}

\section{1.3 Research Objectives}
The objectives of this thesis are:
\begin{enumerate}
  \item Build an end-to-end four-camera 3D perception pipeline for a garage arena.
  \item Consolidate historical scripts and assets in \texttt{Project\_Cam} into one coherent research workflow.
  \item Define and execute reproducible GT protocols for ball and human-joint localization.
  \item Quantify absolute error, bias, precision, and dynamic robustness.
  \item Propose a technically consistent roadmap from perception to launcher control.
\end{enumerate}

\section{1.4 Research Questions and Hypotheses}
\subsection{Table 1.1 Research questions, hypotheses, and validation criteria}

\begin{itemize}
  \item RQ1: Can four fixed USB cameras provide stable millimeter-frame 3D tracking in this garage setup?
\end{itemize}
  H1: Yes, if intrinsics and extrinsics are calibrated in the same deployment state and periodically validated.

\begin{itemize}
  \item RQ2: Which factors dominate reconstruction error under practical operation?
\end{itemize}
  H2: Camera movement, tag visibility imbalance, and multi-view overlap are dominant over detector confidence tuning alone.

\begin{itemize}
  \item RQ3: Can post-processing correction reduce systematic bias to support downstream targeting?
\end{itemize}
  H3: Axis-wise linear correction significantly reduces static bias, but cannot replace geometric recalibration.

\begin{itemize}
  \item RQ4: Is the current system already sufficient for autonomous body-part-targeted launching?
\end{itemize}
  H4: Not yet for autonomous actuation; yes for perception-level guidance and staged HIL integration.

Validation criteria include intrinsics reprojection quality, extrinsics overlay quality, static GT error metrics, dynamic stability metrics, and joint localization consistency.

\section{1.5 Scope}
\subsection{In Scope}
\begin{itemize}
  \item Four-camera capture and synchronization.
  \item ChArUco intrinsics and AprilTag extrinsics calibration.
  \item 3D triangulation of ball and 17-joint human pose.
  \item Unified 3D rendering in garage arena frame.
  \item Ball and joint GT protocols with statistical error reporting.
  \item System design for future launcher integration.
\end{itemize}

\subsection{Out of Scope (Current Thesis Stage)}
\begin{itemize}
  \item Full closed-loop autonomous launcher operation with real actuation.
  \item Certified safety validation under production-risk conditions.
  \item Real-time voice model training and multilingual command robustness.
  \item Advanced biomechanics personalization and long-term athlete adaptation.
\end{itemize}

\section{1.6 Methodological Position}
This thesis follows an engineering research approach rather than demonstration-only implementation. Claims are supported with:
\begin{itemize}
  \item repeatable protocols,
  \item explicit assumptions,
  \item quantitative metrics,
  \item failure analysis,
  \item and reproducible scripts/artifacts.
\end{itemize}

In this sense, the work is not evaluated by one "best demo" video, but by statistically summarized behavior across controlled and stress-test conditions.

\section{1.7 Practical Relevance}
For football training environments where budget or infrastructure constraints limit specialized sensor deployment, a calibrated camera-only approach is attractive. It can define and update target geometry in software, reduce hardware complexity, and support replay-based analysis for both coaching and system diagnostics. The same architecture can also extend to rehabilitation drills, reaction-time training, or human-robot interaction studies.

\section{1.8 Chapter Roadmap}
Chapter 2 reviews relevant literature and identifies the contribution gap. Chapter 3 explains how all major folders in \texttt{Project\_Cam} connect into one coherent system. Chapter 4 formalizes mathematical and architectural foundations. Chapter 5 details implementation and runtime decisions. Chapter 6 describes experimental and GT protocols. Chapter 7 reports measured performance and error analysis. Chapter 8 defines the launcher integration blueprint. Chapter 9 concludes with limitations and future research.


\section{1.9 Research Deliverables by Thesis Stage}
To keep the manuscript aligned with research-based evaluation, the final outputs are organized by maturity level:

Concept outputs (completed):
\begin{itemize}
  \item problem framing from open-loop launching to perception-guided targeting,
  \item hypothesis set and measurable validation criteria,
  \item full repository map connecting early scripts and current integrated workflow.
\end{itemize}

Engineering outputs (completed):
\begin{itemize}
  \item fixed camera role mapping and reproducible recording scripts,
  \item calibrated intrinsics and robust extrinsics workflow,
  \item synchronized four-camera processing from raw clips to 3D trajectories,
  \item arena visualization where camera poses, tags, skeleton, and ball share one frame.
\end{itemize}

Scientific outputs (completed):
\begin{itemize}
  \item static and dynamic GT protocols for ball localization,
  \item joint-touch protocol for human keypoint localization,
  \item quantitative error reports with bias decomposition,
  \item correction model analysis and limits.
\end{itemize}

Translational outputs (partially completed):
\begin{itemize}
  \item architecture for command-driven launcher control,
  \item safety-aware no-fire and confidence gating logic,
  \item launcher integration phase plan with measurable exit criteria.
\end{itemize}

This explicit separation is important during defense because it avoids ambiguity between what has been proven and what remains as development work.

\section{1.10 Expected Thesis Reading Strategy for Committee}
The thesis can be read in two passes:

\begin{enumerate}
  \item Scientific pass: Chapters 1, 2, 4, 6, and 7 establish hypotheses, method, protocols, and evidence.
  \item Engineering pass: Chapters 3, 5, and 8 explain implementation lineage, design trade-offs, and deployment roadmap.
\end{enumerate}

This dual-pass structure is intentional for an ECE committee where some members focus on methodology validity and others focus on system integration viability.

\section{1.11 Success Criteria for the Current Semester}
The current semester success criteria were defined as:
\begin{itemize}
  \item creating a unified world-frame system that includes arena and dynamic entities,
  \item quantifying localization errors rather than relying on visual impression,
  \item preparing control-ready interfaces and safety constraints for future launcher integration.
\end{itemize}

These criteria were met at the perception and validation level. Autonomous launcher firing remains future work and is declared as such throughout the manuscript.
